\section{Síntesis y ampliación}

\subsection{Puntos clave}

\begin{enumerate}
    \item \textbf{AutoGen} toma la conversación como abstracción central, admite la composición recursiva de múltiples Agentes y es adecuado para construir sistemas complejos de colaboración entre Agentes.
    \item \textbf{LlamaIndex} parte de la calidad de los datos y construye asistentes de conocimiento de nivel de producción mediante RAG multimodal y razonamiento Agentic.
    \item Las arquitecturas de múltiples Agentes tienen ventajas notables frente a un solo Agente en la descomposición de tareas, la garantía de seguridad y la capacidad de mantenimiento.
    \item Entre las arquitecturas de Agente restringidas y las de propósito general existe un compromiso entre confiabilidad y flexibilidad.
    \item El despliegue en producción debe considerar la conversión a microservicios, las colas de mensajes y human-in-the-loop.
\end{enumerate}

\subsection{Lecturas adicionales}

\begin{itemize}
    \item Wu et al., ``AutoGen: Enabling Next-Gen LLM Applications via Multi-Agent Conversation,'' 2023.
    \item Liu, ``LlamaIndex: A Data Framework for LLM Applications,'' 2022--2024.
    \item Lewis et al., ``Retrieval-Augmented Generation for Knowledge-Intensive NLP Tasks,'' NeurIPS 2020.
\end{itemize}

\end{document}
